\newpage
\begin{center}
  \textbf{\large ПРИЛОЖЕНИЕ А}
\end{center}
\refstepcounter{chapter}
\addcontentsline{toc}{chapter}{ПРИЛОЖЕНИЕ А}
\label{app:dataset}

\phantomsection
\section*{Структура используемого набора данных}
\addcontentsline{toc}{section}{Структура используемого набора данных}

Используемый набор данных содержит две группы записей, соответствующие различным фазам экспериментального протокола. Первая группа относится к этапу предъявления и запоминания стимулов, вторая - к этапу воспоминания и визуального воображения. В основной задаче реконструкции используются записи второй группы, поскольку именно они соответствуют EEG-сигналам, зарегистрированным во время мысленного воспроизведения изображения.

Данные организованы иерархически: верхний уровень соответствует испытуемому, следующий уровень - отдельной экспериментальной пробе. Внутри каждой пробы хранятся синхронизированные EEG- и EOG-записи, а также разметка целевых блоков. EEG-записи содержат многоканальные сигналы электроэнцефалографии, EOG-записи - данные, связанные с движениями глаз и контролем взгляда.

Разметка содержит список блоков, каждый из которых описывает один целевой фрагмент эксперимента. Используемые смысловые поля разметки приведены в табл. \ref{tab:labels-json-format}.

\begin{table}[!t]
  \centering
  \small
  \setlength{\tabcolsep}{3pt}
  \renewcommand{\arraystretch}{1.15}
  \caption{Основные смысловые поля разметки.}
  \label{tab:labels-json-format}
  \begin{tabularx}{\textwidth}{|L{42mm}|Y|}
    \hline
    Поле разметки & Содержание \\
    \hline
    Список блоков & Перечень целевых блоков внутри экспериментальной пробы. \\
    \hline
    Индекс блока & Порядковый номер блока, с которым связаны EEG- и EOG-записи. \\
    \hline
    Тип стимула & Геометрический паттерн или случайный паттерн. \\
    \hline
    Идентификатор паттерна & Номер паттерна внутри набора стимулов. \\
    \hline
    Целевое изображение & Бинарная матрица размера $6 \times 6$, где значения 0 и 1 задают структуру целевого изображения. \\
    \hline
  \end{tabularx}
\end{table}

Таким образом, каждая EEG-эпоха связывается с целевым изображением через индекс блока и соответствующую запись разметки. Перед обучением моделей целевая матрица изображения векторизуется в последовательность из 36 бинарных значений, что позволяет рассматривать реконструкцию как многовыходную бинарную задачу.

\phantomsection
\section*{Формулы метрик оценки реконструкции}
\addcontentsline{toc}{section}{Формулы метрик оценки реконструкции}
\label{app:metrics}

Пусть для одного примера истинное изображение и реконструкция представлены бинарными векторами $\mathbf{y}, \hat{\mathbf{y}} \in \{0,1\}^{m}$, где $m = 36$. Для вычисления метрик используются следующие величины:

\begin{align}
TP &= \sum_{j=1}^{m} \mathbb{I}(y_j = 1 \land \hat{y}_j = 1), \\
TN &= \sum_{j=1}^{m} \mathbb{I}(y_j = 0 \land \hat{y}_j = 0), \\
FP &= \sum_{j=1}^{m} \mathbb{I}(y_j = 0 \land \hat{y}_j = 1), \\
FN &= \sum_{j=1}^{m} \mathbb{I}(y_j = 1 \land \hat{y}_j = 0).
\end{align}

Здесь $TP$ обозначает число верно восстановленных активных пикселей, $TN$ - число верно восстановленных неактивных пикселей, $FP$ - число ложноположительных активных пикселей, а $FN$ - число пропущенных активных пикселей.

Pixel Accuracy определяется как доля правильно восстановленных пикселей:

\begin{equation}
\operatorname{PixelAccuracy} =
\frac{TP + TN}{TP + TN + FP + FN}.
\end{equation}

Precision, recall и F1-score относятся к стандартным метрикам бинарной классификации~\footcite{10.1016/j.ipm.2009.03.002}:

\begin{align}
\operatorname{Precision} &= \frac{TP}{TP + FP}, \\
\operatorname{Recall} &= \frac{TP}{TP + FN}, \\
\operatorname{F1} &= 2 \cdot
\frac{\operatorname{Precision} \cdot \operatorname{Recall}}
{\operatorname{Precision} + \operatorname{Recall}}.
\end{align}

Метрика Intersection over Union (IoU), также связанная с Jaccard overlap, оценивает совпадение активных областей истинного и восстановленного изображения~\footcite{10.1186/s12880-015-0068-x}:

\begin{equation}
\operatorname{IoU} =
\frac{TP}{TP + FP + FN}.
\end{equation}

Hamming loss показывает долю несовпадающих пикселей между истинным и восстановленным бинарным вектором. Эта метрика основана на идее расстояния Хэмминга~\footcite{10.1002/j.1538-7305.1950.tb00463.x}:

\begin{equation}
\operatorname{HammingLoss} =
\frac{1}{m}\sum_{j=1}^{m}\mathbb{I}(y_j \ne \hat{y}_j).
\end{equation}

При выраженном дисбалансе активных и неактивных пикселей дополнительно может использоваться balanced accuracy:

\begin{equation}
\operatorname{BalancedAccuracy} =
\frac{1}{2}
\left(
\frac{TP}{TP + FN} +
\frac{TN}{TN + FP}
\right).
\end{equation}

Точность полного совпадения является наиболее строгой метрикой: реконструкция
считается правильной только в том случае, если совпали все 36 пикселей:

\begin{equation}
\operatorname{ExactMatch} =
\mathbb{I}(\mathbf{y} = \hat{\mathbf{y}}).
\end{equation}

Для набора данных указанные метрики усредняются по всем примерам тестовой
выборки. В данной работе точность полного совпадения рассматривается только
как вспомогательная метрика, поскольку для малых EEG-датасетов и изображения
из 36 пикселей она может быть чрезмерно жёсткой.

\phantomsection
\section*{Описание признаковых групп}
\addcontentsline{toc}{section}{Описание признаковых групп}
\label{app:feature-groups}

В классических экспериментах EEG-фрагмент сначала приводится к единому
аналитическому представлению: из записи выбирается интервал $[0.5, 15.5)$
с, сигнал ресэмплируется до 125 Гц, а затем признаки вычисляются по полному
15-секундному окну. EOG-каналы в признаковый вектор не включаются.

\begin{table}[!t]
  \centering
  \small
  \setlength{\tabcolsep}{3pt}
  \renewcommand{\arraystretch}{1.15}
  \caption{Группы признаков и их алгоритмическое содержание.}
  \label{tab:feature-implementation-anchors}
  \begin{tabularx}{\textwidth}{|L{34mm}|L{48mm}|Y|}
    \hline
    Группа & Содержание & Выходное представление \\
    \hline
    Временные признаки & Амплитудные статистики, динамика первого и второго
    различий, Hjorth mobility и Hjorth complexity &
    Набор признаков для каждого EEG-канала \\
    \hline
    Спектральные признаки & Абсолютные и относительные мощности в диапазонах
    delta, theta, alpha, beta и low gamma; суммарные спектральные показатели &
    Частотные признаки для каждого EEG-канала \\
    \hline
    Пространственные признаки & OAS-ковариация, корреляционная матрица и
    симметричный матричный логарифм ковариации &
    Матрицы канал--канал \\
    \hline
    Локальные паттерны & LNDP, 1D-LGP и 1D-LBP; 256-bin вероятностные
    гистограммы на канал при $m=8$ &
    Гистограммы локальных кодов для каждого EEG-канала \\
    \hline
  \end{tabularx}
\end{table}

\subsection*{Временные признаки}

Пусть $x_c[t]$ обозначает сигнал канала $c$ внутри анализируемого окна длины
$T$. Для каждого канала вычисляются среднее, дисперсия, стандартное
отклонение, root mean square (RMS), медиана, median absolute deviation (MAD),
размах peak-to-peak, skewness и excess kurtosis. Дополнительно используются
признаки динамики:

\begin{align}
\operatorname{LineLength}_c &=
\frac{1}{T - 1}\sum_{t=1}^{T-1}|x_c[t] - x_c[t-1]|, \\
\operatorname{ZCR}_c &=
\frac{1}{T - 1}\sum_{t=1}^{T-1}
\mathbb{I}\left(\operatorname{sign}(x_c[t]) \ne
\operatorname{sign}(x_c[t-1])\right).
\end{align}

Hjorth-параметры описывают относительную изменчивость сигнала и его
производных. В реализации используются дисперсии первого и второго конечных
различий:

\begin{align}
\operatorname{Mobility}_c &=
\sqrt{\frac{\operatorname{Var}(\Delta x_c)}
{\operatorname{Var}(x_c)}}, \\
\operatorname{Complexity}_c &=
\frac{
\sqrt{\operatorname{Var}(\Delta^2 x_c) /
\operatorname{Var}(\Delta x_c)}
}{
\operatorname{Mobility}_c
}.
\end{align}

Эти признаки относятся к компактным описательным EEG-характеристикам и
используются как интерпретируемая базовая группа для классических моделей
~\footcite{10.1088/1741-2552/aab2f2}.

\subsection*{Спектральные признаки}

Спектральная группа строится по оценке спектральной плотности мощности
(power spectral density, PSD). Для признакового пайплайна PSD вычисляется
через FFT/Welch-процедуру~\footcite{10.1109/TAU.1967.1161901}; дополнительные
спектральные представления, использованные в нейросетевых экспериментах и
проверочных ноутбуках, включают Morlet wavelet transform~\footcite{10.1137/0515056},
Superlet transform~\footcite{10.1038/s41467-020-20539-9} и short-time Fourier
transform~\footcite{gabor1946theory}. Такая группа признаков мотивирована тем,
что визуальное воображение и связанные с ним EEG-эффекты могут проявляться
в частотной структуре сигнала, в частности в альфа-диапазоне
~\footcite{10.1016/j.cub.2020.04.074}.

Для каждого канала вычисляется мощность в диапазонах
delta $(2,4)$ Гц, theta $(4,8)$ Гц, alpha $(8,13)$ Гц,
beta $(13,30)$ Гц и low gamma $(30,40)$ Гц. Если $S_c(f)$ - PSD канала,
то абсолютная мощность диапазона $B=[f_1,f_2)$ определяется как

\begin{equation}
P_{c,B} = \int_{f_1}^{f_2} S_c(f)\,df.
\end{equation}

В дискретной реализации интеграл заменяется суммированием по частотным
bin-ам с учётом перекрытия границ диапазона. Относительная мощность
нормируется на суммарную мощность в диапазоне 2--40 Гц:

\begin{equation}
R_{c,B} =
\frac{P_{c,B}}{\sum_{B'} P_{c,B'}}.
\end{equation}

Также вычисляются доминирующая частота, спектральный центроид и
нормированная спектральная энтропия:

\begin{align}
f^{\ast}_c &= \arg\max_f S_c(f), \\
\operatorname{Centroid}_c &=
\frac{\sum_f f S_c(f)}{\sum_f S_c(f)}, \\
\operatorname{Entropy}_c &=
-\frac{1}{\log N}\sum_{f=1}^{N} p_c(f)\log p_c(f),
\qquad
p_c(f) = \frac{S_c(f)}{\sum_{f'}S_c(f')}.
\end{align}

Минимальная воспроизводимая схема вычисления спектральных признаков состоит
из четырёх шагов: выбрать анализируемый интервал EEG, ресэмплировать сигнал
до 125 Гц, оценить PSD на сетке 2--40 Гц и проинтегрировать мощность по
заданным частотным диапазонам.

Спектральные тензоры для нейросетевых моделей
строятся тем же принципом: сигнал ресэмплируется до 125 Гц, используется
частотная сетка 2--40 Гц с шагом 1 Гц, а выход сохраняет форму
канал--частота для FFT/Welch или канал--частота--время для Morlet, Superlet
и STFT.

В FFT/Welch-ветке используется оконное FFT-представление и последующее
переразбиение PSD на целевую частотную сетку. Morlet-ветка задаёт число
циклов как ограниченную функцию частоты, удаляет краевые отсчёты и усредняет
мощность по временным bin-ам. Superlet-ветка использует adaptive Superlet
orders от 1 до 10 с базовым числом циклов $c_1=3$ и также выполняет
краевую обрезку с bin-агрегацией. STFT-ветка использует Hann-окно длиной
2 с, шаг 32 отсчёта и one-sided PSD с последующим переразбиением частот.

\subsection*{Пространственные признаки}

Пространственная группа описывает совместную структуру EEG-каналов. Для
каждого окна вычисляется устойчивая оценка ковариационной матрицы
$\Sigma \in \mathbb{R}^{C \times C}$ с помощью Oracle Approximating Shrinkage
(OAS). Ковариационные представления широко используются в EEG-BCI и связаны
с римановым подходом к анализу симметричных положительно определённых матриц
~\footcite{10.1080/2326263x.2017.1297192}. При интерпретации межканальных
зависимостей важно учитывать объёмную проводимость, из-за которой измеренные
каналы могут содержать смешанные вклады общих источников~\footcite{10.1002/hbm.20745}.

Из ковариации строится корреляционная матрица:

\begin{equation}
\operatorname{Corr}_{ij} =
\frac{\Sigma_{ij}}{\sqrt{\Sigma_{ii}\Sigma_{jj}}},
\end{equation}

а также симметричный матричный логарифм ковариации. Если
$\Sigma = Q\Lambda Q^\top$ - собственное разложение, то

\begin{equation}
\log(\Sigma) = Q \log(\Lambda) Q^\top,
\end{equation}

где логарифм применяется к собственным значениям. В реализации малые
собственные значения ограничиваются снизу относительно максимального
собственного значения, чтобы избежать численной нестабильности.

\subsection*{Локальные паттерны}

Локально-паттерновые признаки основаны на одномерных расширениях LBP-подхода
для EEG-сигналов: LNDP, 1D-LGP и 1D-LBP~\footcite{10.1016/j.bspc.2017.01.005}.
Для каждого канала рассматриваются скользящие локальные фрагменты длины
$m+1$, где в основной конфигурации $m=8$. Каждый фрагмент преобразуется в
целочисленный код, после чего по кодам строится гистограмма длины
$2^m = 256$. В основной конфигурации гистограмма нормируется до вероятностей.

Для 1D-LBP биты формируются сравнением соседей с центральным отсчётом:

\begin{equation}
b_k = \mathbb{I}(p_k - p_{\mathrm{center}} \ge 0).
\end{equation}

Для 1D-LGP используются абсолютные локальные градиенты относительно центра
и сравнение каждого градиента со средним градиентом внутри фрагмента:

\begin{equation}
b_k =
\mathbb{I}\left(
|p_k - p_{\mathrm{center}}| -
\frac{1}{m}\sum_{r=1}^{m}|p_r - p_{\mathrm{center}}| \ge 0
\right).
\end{equation}

Для LNDP код строится по знакам последовательных локальных различий:

\begin{equation}
b_k = \mathbb{I}(p_{k+1} - p_k \ge 0).
\end{equation}

Полученный набор битов переводится в код

\begin{equation}
\operatorname{code} = \sum_{k=0}^{m-1} b_k 2^k,
\end{equation}

а затем в вероятностную гистограмму:

\begin{equation}
h_j =
\frac{1}{N}\sum_{n=1}^{N}\mathbb{I}(\operatorname{code}_n = j),
\qquad
j = 0,\dots,2^m-1.
\end{equation}

Минимальная схема вычисления локально-паттерновых признаков такова: для
каждого канала сформировать все полные локальные фрагменты длины $m+1$,
преобразовать каждый фрагмент в LNDP-, 1D-LGP- или 1D-LBP-код, затем
построить нормированную гистограмму кодов.

В итоговом признаковом пайплайне блоки разных групп сохраняют свою структуру:
канальные признаки остаются наборами величин для отдельных каналов,
пространственные признаки - матрицами канал--канал, а локальные паттерны -
канальными гистограммами. При обучении конкретной классической модели эти
блоки приводятся к векторному представлению только внутри соответствующего
протокола, настраиваемого только по обучающей части, чтобы масштабирование,
отбор признаков и настройка
гиперпараметров не использовали информацию из тестовой части.

\phantomsection
\section*{Архитектуры нейросетевых моделей}
\addcontentsline{toc}{section}{Архитектуры нейросетевых моделей}
\label{app:deep-architectures}

Нейросетевые эксперименты используют спектральные адаптации архитектур из
семейства ARL EEGModels. Реализация доступна в модуле
\url{https://github.com/Slauva/-master-thesis-code-2024-2026}.
Входной тензор каждой модели имеет вид

\begin{equation}
X \in \mathbb{R}^{B \times P \times C \times W},
\end{equation}

где $B$ - размер батча, $P$ - число спектральных плоскостей, $C$ - число
EEG-электродов, а $W$ - ширина спектральной или временно-спектральной оси.
Для FFT/Welch вход имеет форму $(B,1,C,F)$, а для Morlet, Superlet и STFT -
форму $(B,F,C,T)$, где частотные отсчёты становятся плоскостями входа.
EOG-каналы в эти модели не подаются. Все архитектуры возвращают 36
вещественных логитов, соответствующих пикселям целевого бинарного изображения;
sigmoid-преобразование, порог 0.5 и вычисление метрик выполняются вне
архитектурного модуля.

Важно, что используемые модели не являются численно идентичными переносами
оригинальных TensorFlow/Keras-реализаций. Они сохраняют основные идеи
исходных EEGNet, DeepConvNet и ShallowConvNet~\footcite{10.1088/1741-2552/aace8c,10.1002/hbm.23730},
но адаптированы к спектральным входам: ядра ограничиваются доступной шириной
оси $W$, используется TensorFlow-style SAME padding, а в конце применяется
adaptive global pooling вместо фиксированной flatten-геометрии исходных
моделей. Ограничения max-norm применяются явно после шага оптимизатора.

\begin{table}[!t]
  \centering
  \small
  \setlength{\tabcolsep}{3pt}
  \renewcommand{\arraystretch}{1.15}
  \caption{Нейросетевые архитектуры и их роль в экспериментальном пайплайне.}
  \label{tab:deep-architecture-summary}
  \begin{tabularx}{\textwidth}{|L{34mm}|L{38mm}|Y|}
    \hline
    Архитектура & Статус в работе & Основная структура реализации \\
    \hline
    EEGNet & Основная архитектура полного сравнения &
    Временная свёртка, depthwise-пространственная свёртка по всем электродам,
    separable-блок, average pooling, dropout и линейная голова на 36 выходов. \\
    \hline
    DeepConvNet & Основная архитектура полного сравнения &
    Временная и пространственная свёртки в начальном блоке, затем
    последовательные convolution--batch normalization--ELU--max pooling--dropout
    блоки с 25, 50, 100 и 200 фильтрами и линейная голова. \\
    \hline
    ShallowConvNet & Основная архитектура полного сравнения &
    Временная свёртка, пространственная свёртка, batch normalization,
    квадратичная нелинейность, average pooling, логарифмическая нелинейность,
    dropout и линейная голова. \\
    \hline
    EEGNet-SSVEP & Исследовательский порт, не входит в 12 основных запусков &
    Вариант EEGNet с увеличенным временным ядром и числом фильтров
    ($F_1=96$, $F_2=96$), сохранённый для совместимости с ARL EEGModels. \\
    \hline
    EEGNet-v1 & Исследовательский порт, не входит в 12 основных запусков &
    Более ранняя EEGNet-подобная схема: пространственная свёртка, перестановка
    осей, две двумерные свёртки с batch normalization, ELU и dropout. \\
    \hline
  \end{tabularx}
\end{table}

\subsection*{EEGNet}

EEGNet используется как компактная архитектура, в которой явно разделены
темпоральная, пространственная и separable-свёртки. В спектральной адаптации
первый блок выполняет свёртку только вдоль оси $W$ и формирует $F_1=8$
карт признаков. Затем depthwise-свёртка с ядром $(C,1)$ агрегирует информацию
по всем электродам отдельно для каждой темпоральной карты. При множителе
глубины $D=2$ этот блок создаёт $F_1D=16$ карт признаков. После ELU,
average pooling и dropout применяется separable-блок: depthwise-свёртка
вдоль оси $W$, затем pointwise-свёртка $1\times1$ до $F_2=16$ карт.
Финальная adaptive average pooling-операция сводит пространственно-временные
оси к одному вектору, который подаётся в линейный слой из 36 выходов.

Минимальная схема прямого прохода:

\begin{enumerate}
  \item $X \rightarrow$ SAME temporal convolution $(1, \min(64,W))$.
  \item Batch normalization $\rightarrow$ depthwise spatial convolution $(C,1)$.
  \item ELU $\rightarrow$ average pooling $(1,4)$ $\rightarrow$ dropout.
  \item Depthwise separable convolution $\rightarrow$ ELU $\rightarrow$
  average pooling $(1,8)$ $\rightarrow$ dropout.
  \item Adaptive global pooling $\rightarrow$ linear head $\rightarrow 36$ logits.
\end{enumerate}

\subsection*{DeepConvNet}

DeepConvNet выступает как более глубокая свёрточная базовая архитектура для
EEG-декодирования~\footcite{10.1002/hbm.23730}. В реализации первый блок
разделяет свёртку вдоль спектральной оси и пространственную свёртку по всем
электродам. После этого используются три дополнительные свёрточные ступени.
Число карт признаков возрастает по схеме $25 \rightarrow 50 \rightarrow
100 \rightarrow 200$. Каждый блок включает batch normalization, ELU,
max pooling и dropout. После adaptive global pooling применяется линейная
голова из 36 выходов.

Диаграммно эту архитектуру удобно показывать как вертикальную цепочку блоков:
«спектральная свёртка 25» $\rightarrow$ «пространственная свёртка 25»
$\rightarrow$ «Conv 50» $\rightarrow$ «Conv 100» $\rightarrow$
«Conv 200» $\rightarrow$ «36 логитов». Для всех свёрточных слоёв и линейной
головы в реализации предусмотрена max-norm-проекция.

\subsection*{ShallowConvNet}

ShallowConvNet предназначена для более прямого воспроизведения классической
схемы band-power-декодирования в виде нейросетевого слоя~\footcite{10.1002/hbm.23730}.
Сначала применяется свёртка вдоль оси $W$ с 40 фильтрами, затем
пространственная свёртка с ядром $(C,1)$. После batch normalization
используется квадратичная нелинейность, average pooling и логарифмирование
с нижним численным порогом. Такая последовательность соответствует идее:
выделить фильтрованный сигнал, оценить локальную мощность и подать её в
классификационную голову. В данной задаче классификационная голова заменена
на многовыходный слой из 36 логитов.

\subsection*{Исследовательские порты EEGNet-SSVEP и EEGNet-v1}

Помимо трёх основных архитектур, в коде реализованы EEGNet-SSVEP и EEGNet-v1.
Они использовались как проверочные порты совместимости с ARL EEGModels и не
входят в основную матрицу из 12 полных запусков, где пересекаются три
архитектуры EEGNet, DeepConvNet и ShallowConvNet с четырьмя спектральными
представлениями FFT/Welch, Morlet, Superlet и STFT.

EEGNet-SSVEP наследует структуру EEGNet, но использует более длинное
исходное ядро и больше фильтров. В спектральной адаптации длина ядра также
ограничивается фактической шириной входа $W$. EEGNet-v1 сначала выполняет
пространственную свёртку по электродам, затем переставляет оси и применяет
две двумерные свёртки с batch normalization, ELU и dropout. Для обеих
моделей сохраняется общий контракт: вход $(B,P,C,W)$, выход 36 логитов,
sigmoid и пороговое решение вне модели.

\subsection*{Схемы основных архитектур}

Основные архитектуры полного нейросетевого сравнения показаны на рис.
\ref{fig:eegnet-architecture}--\ref{fig:shallowconvnet-architecture}. Схемы
отражают общий порядок блоков в спектральных адаптациях и не задают
тензорные размеры каждого промежуточного слоя.

\begin{figure}[!t]
  \centering
  \includegraphics[width=150mm]{eegnet_architecture.jpg}
  \caption{Схема спектральной адаптации EEGNet для многовыходной реконструкции.}
  \label{fig:eegnet-architecture}
\end{figure}

\begin{figure}[!t]
  \centering
  \includegraphics[width=150mm]{deepconvnet_architecture.jpg}
  \caption{Схема спектральной адаптации DeepConvNet для многовыходной реконструкции.}
  \label{fig:deepconvnet-architecture}
\end{figure}

\begin{figure}[!t]
  \centering
  \includegraphics[width=150mm]{shallowconvnet_architecture.jpg}
  \caption{Схема спектральной адаптации ShallowConvNet для многовыходной реконструкции.}
  \label{fig:shallowconvnet-architecture}
\end{figure}

EEGNet-SSVEP и EEGNet-v1 в данном приложении описаны текстово как
исследовательские порты совместимости, поскольку они не входят в основную
матрицу полных запусков.
